\documentclass[11pt,a4paper]{amsart}

\usepackage[utf8]{inputenc}
\usepackage[T1]{fontenc}
\usepackage{amsmath,amssymb,amsthm}
\usepackage{mathtools}
\usepackage{enumitem}
\usepackage{booktabs}
\usepackage[margin=1in]{geometry}
\usepackage{xcolor}
\usepackage{hyperref}
\hypersetup{
  colorlinks=true,
  linkcolor=blue!60!black,
  citecolor=blue!60!black,
  urlcolor=blue!60!black
}

% Theorem environments
\theoremstyle{plain}
\newtheorem{theorem}{Theorem}[section]
\newtheorem{lemma}[theorem]{Lemma}
\newtheorem{corollary}[theorem]{Corollary}
\newtheorem{proposition}[theorem]{Proposition}

\theoremstyle{definition}
\newtheorem{definition}[theorem]{Definition}
\newtheorem{remark}[theorem]{Remark}

% Operators
\DeclareMathOperator{\opnorm}{op}
\DeclareMathOperator{\spn}{span}

\title{The Free $\ell^1$ Seminorm on Banach Presheaf Coboundaries}

\author{Jeremy H.\ Carroll}

\date{March 2026}

\begin{document}

\begin{abstract}
We classify seminorms on the 1-cochain space of a Banach presheaf over a finite poset satisfying additivity, faithfulness, and contractive monotonicity. Specifically, let $P$ be a finite partially ordered set and $F$ a Banach presheaf over $P$. The \emph{1-cochain space} is $C^1(P, F) = \bigoplus_{e \in E_{\mathrm{cov}}} F(a_e)$, where the direct sum is taken in the $\ell^1$ sense. We prove that the $\ell^1$ coboundary norm $I(s) = \sum_e \|\delta_e(s)\|$ is, up to a positive scalar, the \textbf{unique} seminorm on $C^1(P, F)$ satisfying these three properties. Adding isomorphism invariance forces the scalar to be uniform across edges. No choice is involved; the norm is generated by the constraints.

\medskip\noindent
In short: \emph{any additive, faithful, contractively monotone diagnostic on Banach presheaf coboundaries must be a positive scalar multiple of the $\ell^1$ norm.}
\end{abstract}

\maketitle

\section{Setup}

\subsection{Banach Presheaves over Finite Posets}

Let $(P, \leq)$ be a finite partially ordered set. The \emph{covering relations} are the edges of the Hasse diagram:
\[
E_{\mathrm{cov}}(P) = \{e = (a, b) : a < b,\; \nexists\, c \text{ with } a < c < b\}.
\]

A \emph{Banach presheaf} $F$ over $P$ assigns:
\begin{itemize}[nosep]
  \item To each $x \in P$, a Banach space $(F(x), \|\cdot\|_{F(x)})$.
  \item To each $e = (a, b) \in E_{\mathrm{cov}}$, a bounded linear \emph{restriction map} $r_{b,a}: F(b) \to F(a)$ with $\|r_{b,a}\|_{\opnorm} \leq 1$.
\end{itemize}

A \emph{section} of $F$ is a family $s = (s(x))_{x \in P}$ with $s(x) \in F(x)$.

\subsection{The 1-Cochain Space}

For each edge $e = (a, b)$, the \emph{coboundary} is:
\[
\delta_e(s) = r_{b,a}(s(b)) - s(a) \in F(a).
\]

The \emph{1-cochain space} is the Banach direct sum:
\[
C^1(P, F) = \bigoplus_{e \in E_{\mathrm{cov}}}^{\ell^1} F(a_e)
\]
with norm $\|(\delta_e)_e\|_1 = \sum_e \|\delta_e\|_{F(a_e)}$. This is isometrically isomorphic to the space of coboundary vectors.

The \emph{$\ell^1$ coboundary norm} of a section $s$ is:
\[
I(s) = \sum_{e \in E_{\mathrm{cov}}} \|\delta_e(s)\|_{F(a_e)} = \|(\delta_e(s))_e\|_1.
\]

\subsection{Seminorms on 1-Cochains}

\begin{definition}\label{def:seminorm}
A \emph{seminorm on 1-cochains} is a functional $\nu: C^1(P, F) \to \mathbb{R}_{\geq 0}$.
\end{definition}

We consider four properties:

\medskip
\noindent\textbf{(P1) Edge-wise additivity.} $\nu$ decomposes as
\[
\nu((\delta_e)_e) = \sum_{e \in E_{\mathrm{cov}}} \nu_e(\delta_e)
\]
where each $\nu_e: F(a_e) \to \mathbb{R}_{\geq 0}$ is itself a seminorm on $F(a_e)$.

\medskip
\noindent\textbf{(P2) Faithfulness.} For each $e$: $\nu_e(\delta) = 0 \iff \delta = 0$.

\medskip
\noindent\textbf{(P3) Contractive monotonicity.} For every bounded linear map $\varphi: V \to W$ between Banach spaces, and every presheaf morphism induced stalkwise by $\varphi$:
\[
\nu_e(\varphi(\delta)) \leq \|\varphi\|_{\opnorm} \cdot \nu_e(\delta).
\]

\medskip
\noindent\textbf{(P4) Isomorphism invariance.} For any poset isomorphism $\pi: \tilde{P} \to P$:
\[
\nu_{\tilde{P}}(\pi^*(\delta_e)_e) = \nu_P((\delta_e)_e).
\]

\section{The Classification Theorem}

We first isolate two structural lemmas.

\begin{lemma}[1-Dimensional Compatibility]\label{lem:1d}
Let $\nu_e$ satisfy \textup{(P3)}. Then for any 1-dimensional Banach subspace $\spn(\delta) \subset V$, the restriction $\nu_e|_{\spn(\delta)}$ is well-defined and satisfies \textup{(P1)--(P3)} with respect to the induced Banach norm on $\spn(\delta) \cong \mathbb{R}$. Specifically: any isometric embedding $\iota: \mathbb{R} \hookrightarrow V$ induces a seminorm $\nu_e^{\mathbb{R}} = \nu_e \circ \iota$ on $\mathbb{R}$ satisfying \textup{(P2)--(P3)}, and any contractive surjection $f: V \to \mathbb{R}$ induces $\nu_e^{\mathbb{R}}(f(\delta)) \leq \|f\|_{\opnorm} \cdot \nu_e(\delta)$.
\end{lemma}

\begin{proof}
Direct application of (P3) to $\iota$ and $f$ respectively. Both are bounded linear maps between Banach spaces, so (P3) applies without additional hypothesis.
\end{proof}

\begin{lemma}[Direction Independence]\label{lem:direction}
If $\nu_e$ satisfies \textup{(P3)}, then for any $v, w \in F(a_e)$ with $\|v\| = \|w\| \neq 0$, we have $\nu_e(v) = \nu_e(w)$.
\end{lemma}

\begin{proof}
By the Hahn--Banach theorem, there exists $f_v \in F(a_e)^*$ with $\|f_v\|_{\opnorm} = 1$ and $f_v(v) = \|v\|$. Define $\Phi: F(a_e) \to F(a_e)$ by $\Phi(x) = f_v(x) \frac{w}{\|w\|}$. Then $\Phi(v) = w$ and $\|\Phi\|_{\opnorm} = 1$, so $\nu_e(w) \leq \nu_e(v)$ by (P3). By symmetry, constructing $\Psi(x) = f_w(x) \frac{v}{\|v\|}$ gives $\nu_e(v) \leq \nu_e(w)$, forcing equality.
\end{proof}

\begin{theorem}[Free Seminorm Classification]\label{thm:main}
Let $\nu$ be a seminorm on $C^1(P, F)$ satisfying \textup{(P1)}, \textup{(P2)}, and \textup{(P3)}. Then for each edge $e$, there exists $w_e > 0$ such that
\[
\nu_e(\delta) = w_e \cdot \|\delta\|_{F(a_e)} \quad \text{for all } \delta \in F(a_e).
\]
If additionally \textup{(P4)} holds, then $w_e = \alpha$ is independent of $e$, and
\[
\nu((\delta_e)_e) = \alpha \sum_{e} \|\delta_e\|_{F(a_e)} = \alpha \cdot I(s).
\]
\end{theorem}

In other words: the $\ell^1$ coboundary norm is the unique (up to positive scalar) seminorm satisfying all four properties.

\section{Proof}

The proof proceeds in four stages: direction independence, factoring through norms, forcing linearity, and forcing uniformity.

\subsection*{Stage 1: Direction Independence}

This is Lemma~\ref{lem:direction} above.

\subsection*{Stage 2: Each $\nu_e$ depends only on $\|\delta\|$}

Fix an edge $e$ and let $V = F(a_e)$. Fix $\delta \in V$ with $\|\delta\| = r > 0$.

\medskip
\noindent\textbf{Upper bound.} By the Hahn--Banach theorem, there exists $f \in V^*$ with $\|f\|_{\opnorm} = 1$ and $f(\delta) = r$. View $f$ as a bounded linear map $f: V \to \mathbb{R}$ with $\|f\|_{\opnorm} = 1$. By Lemma~\ref{lem:1d} and (P3):
\[
\nu_e^{\mathbb{R}}(f(\delta)) \leq \|f\|_{\opnorm} \cdot \nu_e(\delta) = \nu_e(\delta)
\]
where $\nu_e^{\mathbb{R}} := \nu_e \circ \iota$ for any unit-norm isometric embedding $\iota: \mathbb{R} \hookrightarrow V$. By Lemma~\ref{lem:direction}, two embeddings mapping $1$ to vectors of equal norm produce identical values of $\nu_e$, so $\nu_e^{\mathbb{R}}$ is well defined. Since $f(\delta) = r$, we get $\nu_e^{\mathbb{R}}(r) \leq \nu_e(\delta)$.

\medskip
\noindent\textbf{Lower bound.} Define $\iota: \mathbb{R} \to V$ by $\iota(t) = t \cdot \delta / \|\delta\|$. Then $\|\iota\|_{\opnorm} = 1$ (it maps $t$ to a vector of norm $|t|$). By (P3) applied to $\iota$:
\[
\nu_e(\iota(r)) \leq \|\iota\|_{\opnorm} \cdot \nu_e^{\mathbb{R}}(r) = \nu_e^{\mathbb{R}}(r).
\]
Since $\iota(r) = \delta$, we have $\nu_e(\delta) \leq \nu_e^{\mathbb{R}}(r)$.

\medskip
Combining: $\nu_e(\delta) = \nu_e^{\mathbb{R}}(\|\delta\|)$.

For $\delta = 0$: $\nu_e(0) = 0$ by (P2). Therefore $\nu_e(\delta) = h_e(\|\delta\|)$ where $h_e: \mathbb{R}_{\geq 0} \to \mathbb{R}_{\geq 0}$ with $h_e(0) = 0$ and $h_e(r) > 0$ for $r > 0$.\qed\ (Stage 2)

\subsection*{Stage 3: $h_e$ is linear}

Apply (P3) to the scaling maps $v \mapsto cv$ and $v \mapsto v/c$ for $c > 0$, with operator norms $c$ and $1/c$ respectively. These yield:
\[
h_e(cr) \leq c \cdot h_e(r) \quad \text{and} \quad c \cdot h_e(r) \leq h_e(cr).
\]

Together: $h_e(c \cdot r) = c \cdot h_e(r)$ for all $c > 0$ and $r \geq 0$.

This is exact 1-homogeneity. Setting $r = 1$:
\[
h_e(c) = c \cdot h_e(1) \quad \text{for all } c > 0.
\]

Combined with $h_e(0) = 0$ (faithfulness), this gives $h_e(d) = w_e \cdot d$ on all of $\mathbb{R}_{\geq 0}$, where $w_e = h_e(1) > 0$ (positive by (P2)).\qed\ (Stage 3)

\subsection*{Stage 4: Isomorphism invariance forces uniform weights}

Let $e = (a, b)$ be any covering edge of a finite poset $P$. Define the \emph{single-edge subposet} $P_e = \{a < b\}$ with the restricted presheaf $F_e = F|_{P_e}$, and let $\iota_e: P_e \hookrightarrow P$ denote the inclusion functor.

By (P1), $\nu$ decomposes as $\nu((\delta_e)_e) = \sum_{e'} \nu_{e'}(\delta_{e'})$. Each summand $\nu_{e'}$ depends only on $\delta_{e'} \in F(a_{e'})$ and is determined independently of the other summands. In particular, $\nu_e(\delta_e)$ depends only on the coboundary $\delta_e(s) = r_{b,a}(s(b)) - s(a)$, which is computed entirely from the restricted presheaf data $(F(a), F(b), r_{b,a})$ --- that is, from $F_e$ alone. The value $\nu_e(\delta_e)$ is therefore independent of the ambient poset $P$ in which $e$ is embedded: it is determined by the single-edge presheaf $(P_e, F_e)$.

By Stages 1--3, $\nu_e(\delta_e) = w_e \|\delta_e\|$ for some $w_e > 0$. Since $w_e$ depends only on the isomorphism class of the single-edge presheaf $(P_e, F_e)$, it remains to show all such presheaves yield the same weight.

Let $e_1 = (a_1, b_1)$ and $e_2 = (a_2, b_2)$ be any two covering edges (in the same or different posets) with the same stalk spaces $F(a_{e_1}) \cong F(a_{e_2})$ and $F(b_{e_1}) \cong F(b_{e_2})$. The single-edge posets $P_{e_1} = \{a_1 < b_1\}$ and $P_{e_2} = \{a_2 < b_2\}$ admit a unique order-preserving bijection $\pi: P_{e_1} \to P_{e_2}$. By (P4), $\nu_{P_{e_1}}(\pi^*(\delta_{e_2})) = \nu_{P_{e_2}}((\delta_{e_2}))$, which gives $w_{e_1} = w_{e_2}$.

Thus all covering edges share the same weight $w$.

Setting $\alpha = w$: $\nu = \alpha \cdot I$.\qed\ (Theorem~\ref{thm:main})

\section{Necessity of Each Property}

\begin{center}
\begin{tabular}{@{}lll@{}}
\toprule
Dropped property & What happens & Example \\
\midrule
(P1) Additivity & $\ell^\infty$-type norms admitted & $\nu = \max_e \|\delta_e\|$ \\
(P2) Faithfulness & Trivial seminorm admitted & $\nu \equiv 0$ \\
(P3) Contractivity & Nonlinear seminorms admitted & $\nu = \sum_e \|\delta_e\|^2$ \\
(P4) Invariance & Non-uniform weights admitted & $\nu = \sum_e w_e \|\delta_e\|$ \\
\bottomrule
\end{tabular}
\end{center}

Each property is independently necessary. Together they are sufficient.

\begin{remark}
Properties \textup{(P1)--(P3)} alone force the $\ell^1$-type structure $\nu = \sum_e w_e \|\delta_e\|$ with possibly non-uniform positive weights. Property \textup{(P4)} collapses the weights to a single scalar.
\end{remark}

\section{The Universal Property (Categorical Formulation)}

\begin{definition}\label{def:afc}
Let $\mathbf{AFC}$ (Additive Faithful Contractive) denote the category whose:
\begin{itemize}[nosep]
  \item \textbf{Objects} are seminorms $\nu$ on $C^1(P, F)$ satisfying \textup{(P1)--(P3)} for all finite posets $P$ and Banach presheaves $F$.
  \item \textbf{Morphisms} $\eta: \nu \to \mu$ are given by families of positive reals $\{\lambda_e\}$ with $\mu_e = \lambda_e \cdot \nu_e$ for all $e$.
\end{itemize}
This category is not necessarily small as it spans all Banach presheaves and posets. Furthermore, it is a thin category with respect to any pair of objects, as the morphisms are simply determined by structural scalars relating the seminorms.
\end{definition}

\begin{corollary}[Structure of AFC]\label{cor:afc}
Every object in $\mathbf{AFC}$ has the form $\nu = \sum_e w_e \|\cdot\|_{F(a_e)}$ for some family $\{w_e > 0\}$. In particular, every object in $\mathbf{AFC}$ is a weighted $\ell^1$ norm.
\end{corollary}

\begin{corollary}[Universal Property]\label{cor:universal}
Let $\mathbf{AFC}_4 \subset \mathbf{AFC}$ denote the full subcategory of objects also satisfying \textup{(P4)}. Then:
\begin{enumerate}[label=(\roman*),nosep]
  \item Every object in $\mathbf{AFC}_4$ equals $\alpha \cdot I$ for some $\alpha > 0$.
  \item $I$ is terminal in $\mathbf{AFC}_4$ up to positive scalar normalization: for each $\nu = \alpha \cdot I$, the unique morphism $\nu \to I$ is the global scalar $1/\alpha$.
  \item The universal property diagram commutes: for any $\nu \in \mathbf{AFC}_4$ and any Banach presheaf morphism $\varphi: F \to G$, the square
  \[
  \begin{array}{ccc}
    C^1(P, F) & \xrightarrow{\varphi_*} & C^1(P, G) \\
    \downarrow{\nu} & & \downarrow{\nu} \\
    \mathbb{R}_{\geq 0} & = & \mathbb{R}_{\geq 0}
  \end{array}
  \]
  satisfies $\nu(\varphi_*(\delta)) \leq \|\varphi\|_{\opnorm} \cdot \nu(\delta)$, and $I$ is the unique (up to scalar) seminorm for which this holds.
\end{enumerate}
\end{corollary}

\section{Discussion}

\subsection{What This Paper Proves}

The $\ell^1$ coboundary norm is not selected from a family of candidates. It is \emph{generated} by three structural constraints (additivity, faithfulness, contractivity) on the space of 1-cochains of Banach presheaves over finite posets. Adding relabeling invariance eliminates the only remaining freedom (edge weights). The result is purely functional-analytic --- it requires no physics, no dynamics, and no topology beyond the Hasse diagram of a finite poset.

\subsection{What This Paper Does Not Prove}

\begin{itemize}[nosep]
  \item We do not prove anything about dynamics, time evolution, or physical systems.
  \item We do not extend to infinite posets (where compactness arguments would be needed).
  \item We do not claim the four properties are the ``right'' axioms. We prove that \emph{if} they hold, the norm is forced.
\end{itemize}

\subsection{Relation to Other Work}

The Hahn--Banach factoring argument used in Stage~2 is analogous to classical norm classification techniques in Banach space theory, where linear functionals are used to reduce geometric constraints to one-dimensional analysis (see Megginson~\cite{megginson} for Banach geometric tools, and Pestov~\cite{pestov} for related uniqueness results for norms). The key structural observation is that contractive monotonicity, combined with Hahn--Banach, forces exact 1-homogeneity --- eliminating all nonlinear seminorms without any regularity assumption. This mechanism parallels the uniqueness proof for the operator norm, where submultiplicativity and compatibility with vector norms force norm-dependence via analogous Hahn--Banach arguments. The scaling argument in Stage~3 is elementary but, to our knowledge, has not been applied in this presheaf context (though related categorical sheaf structures appear in Godement~\cite{godement}).

The result is a presheaf version of a classical principle in Banach space theory: the characterization of the $\ell^1$ norm as the unique additive cross-norm on a free Banach sum. The 1-cochain space $C^1(P, F) = \bigoplus_e F(a_e)$ is the \emph{free additive Banach object generated by the edges} of the poset. When a measurement on such a free object must satisfy additivity across components and compatibility (via contractive monotonicity) with the norms on each component, $\ell^1$ is the only geometry that survives. The $\ell^2$ norm fails because $\|(v, w)\|_2 \neq \|v\| + \|w\|$; the $\ell^\infty$ norm fails because $\max(\|v\|, \|w\|)$ discards summation. This is why $\ell^1$ appears so inevitably: it is the \emph{universal additive geometry}.

We note that the coboundary defects $\delta_e(s) = r_{b,a}(s(b)) - s(a)$ play a role analogous to curvature in gauge theories: they measure the failure of local data to glue consistently across the covering relations of the poset. In gauge theory, curvature measures the non-commutativity of parallel transport around loops; here, the coboundary measures the incompatibility of section values across edges. The distinction is that in gauge theory the metric on curvature is a free choice, whereas our classification theorem shows that the defect magnitude is uniquely $\ell^1$ --- the geometry forces the norm.

\subsection{Relation to the Main SOT Paper}

This paper serves as a companion to the Single Obstruction Theory. Together, the two papers form a classical three-layer structure:

\begin{enumerate}[nosep]
  \item \textbf{Existence} (SOT Theorem~3.1): Nonlinear projection introduces curvature --- a differential-geometric obstruction measured by the second-order term $D^2\Pi(Am, v)$, which acts as a second fundamental form.
  \item \textbf{Uniqueness} (this paper, Theorem~\ref{thm:main}): The obstruction must be measured by the $\ell^1$ coboundary norm. No other norm is compatible with the structural constraints.
  \item \textbf{Persistence} (SOT Theorem~5.1): The obstruction density $C(t) \geq c > 0$ persists at finite density for physical time intervals independent of system size.
\end{enumerate}

The main SOT paper asserts $\ell^1$ rigidity as a structural input; this paper proves \emph{why} $\ell^1$ is forced. The separation is deliberate: the classification theorem is a self-contained result in functional analysis that neither requires nor implies any dynamical or physical interpretation. It answers the question: \emph{why $\ell^1$ and not $\ell^2$?} Because the axioms force it.

\begin{thebibliography}{5}

\bibitem{pestov}
V.~Pestov, ``Dynamics of infinite-dimensional groups and Ramsey-type phenomena,'' \emph{Publica\c{c}\~{o}es Matem\'{a}ticas do IMPA}, 2006.

\bibitem{maclane}
S.~Mac~Lane, \emph{Categories for the Working Mathematician}, Springer, 1971.

\bibitem{megginson}
R.~E.~Megginson, \emph{An Introduction to Banach Space Theory}, Springer, 1998.

\bibitem{bourbaki}
N.~Bourbaki, \emph{Topological Vector Spaces}, Springer, 1987.

\bibitem{godement}
R.~Godement, \emph{Topologie Alg\'{e}brique et Th\'{e}orie des Faisceaux}, Hermann, 1958.

\end{thebibliography}

\end{document}
